\documentclass[12pt]{article}
\usepackage[UTF8]{inputenc}
\usepackage{xeCJK}
\setCJKmainfont{Sarabun}
\usepackage{enumitem}
\usepackage{geometry}
\geometry{a4paper, margin=2.5cm}
\usepackage{setspace}
\onehalfspacing

\begin{document}

\begin{center}
{\Large\textbf{สเปกตรัมและสี}}\\[6pt]
{\normalsize วิชาฟิสิกส์ ม.ปลาย บทที่ 24}\\[4pt]
{\footnotesize ทั้งหมด 25 ข้อ}
\end{center}
\hrule
\bigskip

\section*{คำถาม in class}
\begin{enumerate}[label=\arabic*., itemsep=0.5\baselineskip, topsep=0.5\baselineskip, leftmargin=*]
\item แสงขาวเมื่อผ่านปริซึมจะแยกออกเป็นสีต่างๆ เรียกปรากฏการณ์นี้ว่าอะไร?
\begin{enumerate}[label=\alph*)., itemsep=0.3\baselineskip, topsep=0.3\baselineskip]
\item ก\ 
\begin{minipage}[t]{0.9\linewidth}
\raggedright การเลี้ยวเบน
\end{minipage}
\item ข\ 
\begin{minipage}[t]{0.9\linewidth}
\raggedright การสะท้อน
\end{minipage}
\item ค\ 
\begin{minipage}[t]{0.9\linewidth}
\raggedright การกระจายของแสง
\end{minipage}
\item ง\ 
\begin{minipage}[t]{0.9\linewidth}
\raggedright การเคลื่อนที่เชิงคลื่น
\end{minipage}
\end{enumerate}
\vspace{-0.3\baselineskip}
\begin{small}
\textit{คำตอบ: ค}
\end{small}

\item จากสเปกตรัมของแสงขาว สีใดมีความยาวคลื่นสั้นที่สุด?
\begin{enumerate}[label=\alph*)., itemsep=0.3\baselineskip, topsep=0.3\baselineskip]
\item ก\ 
\begin{minipage}[t]{0.9\linewidth}
\raggedright แดง
\end{minipage}
\item ข\ 
\begin{minipage}[t]{0.9\linewidth}
\raggedright เขียว
\end{minipage}
\item ค\ 
\begin{minipage}[t]{0.9\linewidth}
\raggedright เหลือง
\end{minipage}
\item ง\ 
\begin{minipage}[t]{0.9\linewidth}
\raggedright ม่วง
\end{minipage}
\end{enumerate}
\vspace{-0.3\baselineskip}
\begin{small}
\textit{คำตอบ: ง}
\end{small}

\item รังสีใดต่อไปนี้มีความถี่สูงที่สุด?
\begin{enumerate}[label=\alph*)., itemsep=0.3\baselineskip, topsep=0.3\baselineskip]
\item ก\ 
\begin{minipage}[t]{0.9\linewidth}
\raggedright รังสีอินฟราเรด
\end{minipage}
\item ข\ 
\begin{minipage}[t]{0.9\linewidth}
\raggedright รังสีอัลตราไวโอเลต
\end{minipage}
\item ค\ 
\begin{minipage}[t]{0.9\linewidth}
\raggedright รังสีเอกซ์
\end{minipage}
\item ง\ 
\begin{minipage}[t]{0.9\linewidth}
\raggedright รังสีแกมมา
\end{minipage}
\end{enumerate}
\vspace{-0.3\baselineskip}
\begin{small}
\textit{คำตอบ: ง}
\end{small}

\item รังสีอินฟราเรด (Infrared) มีช่วงความยาวคลื่นประมาณเท่าใด?
\begin{enumerate}[label=\alph*)., itemsep=0.3\baselineskip, topsep=0.3\baselineskip]
\item ก\ 
\begin{minipage}[t]{0.9\linewidth}
\raggedright 10⁻⁸ – 10⁻⁶ เมตร
\end{minipage}
\item ข\ 
\begin{minipage}[t]{0.9\linewidth}
\raggedright 10⁻⁶ – 10⁻³ เมตร
\end{minipage}
\item ค\ 
\begin{minipage}[t]{0.9\linewidth}
\raggedright 10⁻³ – 10⁰ เมตร
\end{minipage}
\item ง\ 
\begin{minipage}[t]{0.9\linewidth}
\raggedright 10⁰ – 10³ เมตร
\end{minipage}
\end{enumerate}
\vspace{-0.3\baselineskip}
\begin{small}
\textit{คำตอบ: ข}
\end{small}

\item ตามองเห็นแสงสีต่างๆ ได้เพราะเซลล์ประสาทใดในจอประสาทตา?
\begin{enumerate}[label=\alph*)., itemsep=0.3\baselineskip, topsep=0.3\baselineskip]
\item ก\ 
\begin{minipage}[t]{0.9\linewidth}
\raggedright เซลล์แท่ง
\end{minipage}
\item ข\ 
\begin{minipage}[t]{0.9\linewidth}
\raggedright เซลล์กรวย
\end{minipage}
\item ค\ 
\begin{minipage}[t]{0.9\linewidth}
\raggedright เซลล์ปมประสาท
\end{minipage}
\item ง\ 
\begin{minipage}[t]{0.9\linewidth}
\raggedright เซลล์ไมอีลิน
\end{minipage}
\end{enumerate}
\vspace{-0.3\baselineskip}
\begin{small}
\textit{คำตอบ: ข}
\end{small}

\item การผสมสีแบบ RGB ใช้หลักการผสมแสงสีผสมกัน ถ้าผสมแสงแดงกับแสงเขียวจะได้สีใด?
\begin{enumerate}[label=\alph*)., itemsep=0.3\baselineskip, topsep=0.3\baselineskip]
\item ก\ 
\begin{minipage}[t]{0.9\linewidth}
\raggedright แสงขาว
\end{minipage}
\item ข\ 
\begin{minipage}[t]{0.9\linewidth}
\raggedright แสงเหลือง
\end{minipage}
\item ค\ 
\begin{minipage}[t]{0.9\linewidth}
\raggedright แสงฟ้า
\end{minipage}
\item ง\ 
\begin{minipage}[t]{0.9\linewidth}
\raggedright แสงม่วง
\end{minipage}
\end{enumerate}
\vspace{-0.3\baselineskip}
\begin{small}
\textit{คำตอบ: ข}
\end{small}

\item การผสมสีแบบ CMYK ใช้ในอุปกรณ์พิมพ์ ถ้าต้องการพิมพ์สีแดงต้องผสมสีใด?
\begin{enumerate}[label=\alph*)., itemsep=0.3\baselineskip, topsep=0.3\baselineskip]
\item ก\ 
\begin{minipage}[t]{0.9\linewidth}
\raggedright ฟ้า + เหลือง
\end{minipage}
\item ข\ 
\begin{minipage}[t]{0.9\linewidth}
\raggedright แดงเด่น + เหลือง
\end{minipage}
\item ค\ 
\begin{minipage}[t]{0.9\linewidth}
\raggedright ฟ้า + ม่วงแดง
\end{minipage}
\item ง\ 
\begin{minipage}[t]{0.9\linewidth}
\raggedright เขียว + ม่วงแดง
\end{minipage}
\end{enumerate}
\vspace{-0.3\baselineskip}
\begin{small}
\textit{คำตอบ: ข}
\end{small}

\item กฎการแผ่รังสีของวีน์ (Wien's Law) บอกว่าความยาวคลื่นที่ความเข้มสูงสุดแปรผกผันกับอะไร?
\begin{enumerate}[label=\alph*)., itemsep=0.3\baselineskip, topsep=0.3\baselineskip]
\item ก\ 
\begin{minipage}[t]{0.9\linewidth}
\raggedright มวล
\end{minipage}
\item ข\ 
\begin{minipage}[t]{0.9\linewidth}
\raggedright ความดัน
\end{minipage}
\item ค\ 
\begin{minipage}[t]{0.9\linewidth}
\raggedright อุณหภูมิ
\end{minipage}
\item ง\ 
\begin{minipage}[t]{0.9\linewidth}
\raggedright ความเร็ว
\end{minipage}
\end{enumerate}
\vspace{-0.3\baselineskip}
\begin{small}
\textit{คำตอบ: ค}
\end{small}

\item ช่วงความยาวคลื่นของแสงที่ตามองเห็นได้ (visible light) อยู่ในช่วงใด?
\begin{enumerate}[label=\alph*)., itemsep=0.3\baselineskip, topsep=0.3\baselineskip]
\item ก\ 
\begin{minipage}[t]{0.9\linewidth}
\raggedright 1 – 10 nm
\end{minipage}
\item ข\ 
\begin{minipage}[t]{0.9\linewidth}
\raggedright 10 – 100 nm
\end{minipage}
\item ค\ 
\begin{minipage}[t]{0.9\linewidth}
\raggedright 400 – 700 nm
\end{minipage}
\item ง\ 
\begin{minipage}[t]{0.9\linewidth}
\raggedright 700 – 1000 nm
\end{minipage}
\end{enumerate}
\vspace{-0.3\baselineskip}
\begin{small}
\textit{คำตอบ: ค}
\end{small}

\item รังสีอัลตราไวโอเลต (UV) มีความยาวคลื่นอยู่ระหว่างขอบเขตใด?
\begin{enumerate}[label=\alph*)., itemsep=0.3\baselineskip, topsep=0.3\baselineskip]
\item ก\ 
\begin{minipage}[t]{0.9\linewidth}
\raggedright 1 – 400 nm
\end{minipage}
\item ข\ 
\begin{minipage}[t]{0.9\linewidth}
\raggedright 400 – 700 nm
\end{minipage}
\item ค\ 
\begin{minipage}[t]{0.9\linewidth}
\raggedright 700 – 1000 nm
\end{minipage}
\item ง\ 
\begin{minipage}[t]{0.9\linewidth}
\raggedright 1 – 10 nm
\end{minipage}
\end{enumerate}
\vspace{-0.3\baselineskip}
\begin{small}
\textit{คำตอบ: ก}
\end{small}

\item ดวงอาทิตย์มีอุณหภูมิพื้นผิวประมาณ 5800 K ความยาวคลื่นที่ความเข้มสูงสุดของดวงอาทิตย์อยู่ที่ประมาณเท่าใด?
\begin{enumerate}[label=\alph*)., itemsep=0.3\baselineskip, topsep=0.3\baselineskip]
\item ก\ 
\begin{minipage}[t]{0.9\linewidth}
\raggedright 500 nm
\end{minipage}
\item ข\ 
\begin{minipage}[t]{0.9\linewidth}
\raggedright 1000 nm
\end{minipage}
\item ค\ 
\begin{minipage}[t]{0.9\linewidth}
\raggedright 1500 nm
\end{minipage}
\item ง\ 
\begin{minipage}[t]{0.9\linewidth}
\raggedright 2000 nm
\end{minipage}
\end{enumerate}
\vspace{-0.3\baselineskip}
\begin{small}
\textit{คำตอบ: ก}
\end{small}

\item วัตถุสีแดงเมื่อส่องด้วยแสงขาวจะเห็นเป็นสีแดงเพราะเหตุผลใด?
\begin{enumerate}[label=\alph*)., itemsep=0.3\baselineskip, topsep=0.3\baselineskip]
\item ก\ 
\begin{minipage}[t]{0.9\linewidth}
\raggedright ดูดกลืนแสงแดง สะท้อนสีอื่น
\end{minipage}
\item ข\ 
\begin{minipage}[t]{0.9\linewidth}
\raggedright ดูดกลืนสีอื่น สะท้อนแสงแดง
\end{minipage}
\item ค\ 
\begin{minipage}[t]{0.9\linewidth}
\raggedright กระจายแสงแดงออกไป
\end{minipage}
\item ง\ 
\begin{minipage}[t]{0.9\linewidth}
\raggedright เปล่งแสงแดงออกมา
\end{minipage}
\end{enumerate}
\vspace{-0.3\baselineskip}
\begin{small}
\textit{คำตอบ: ข}
\end{small}

\item ในสเปกตรัมของแสงขาว สีใดมีความถี่สูงที่สุด?
\begin{enumerate}[label=\alph*)., itemsep=0.3\baselineskip, topsep=0.3\baselineskip]
\item ก\ 
\begin{minipage}[t]{0.9\linewidth}
\raggedright แดง
\end{minipage}
\item ข\ 
\begin{minipage}[t]{0.9\linewidth}
\raggedright เขียว
\end{minipage}
\item ค\ 
\begin{minipage}[t]{0.9\linewidth}
\raggedright เหลือง
\end{minipage}
\item ง\ 
\begin{minipage}[t]{0.9\linewidth}
\raggedright ม่วง
\end{minipage}
\end{enumerate}
\vspace{-0.3\baselineskip}
\begin{small}
\textit{คำตอบ: ง}
\end{small}

\item ถ้าผสมแสงสีทั้งสาม RGB เต็มความเข้มจะได้สีใด?
\begin{enumerate}[label=\alph*)., itemsep=0.3\baselineskip, topsep=0.3\baselineskip]
\item ก\ 
\begin{minipage}[t]{0.9\linewidth}
\raggedright ดำ
\end{minipage}
\item ข\ 
\begin{minipage}[t]{0.9\linewidth}
\raggedright ขาว
\end{minipage}
\item ค\ 
\begin{minipage}[t]{0.9\linewidth}
\raggedright เหลือง
\end{minipage}
\item ง\ 
\begin{minipage}[t]{0.9\linewidth}
\raggedright เทา
\end{minipage}
\end{enumerate}
\vspace{-0.3\baselineskip}
\begin{small}
\textit{คำตอบ: ข}
\end{small}

\item รังสีเอกซ์ (X-ray) มีความยาวคลื่นอยู่ในช่วงใดโดยประมาณ?
\begin{enumerate}[label=\alph*)., itemsep=0.3\baselineskip, topsep=0.3\baselineskip]
\item ก\ 
\begin{minipage}[t]{0.9\linewidth}
\raggedright 0.01 – 10 nm
\end{minipage}
\item ข\ 
\begin{minipage}[t]{0.9\linewidth}
\raggedright 10 – 400 nm
\end{minipage}
\item ค\ 
\begin{minipage}[t]{0.9\linewidth}
\raggedright 400 – 700 nm
\end{minipage}
\item ง\ 
\begin{minipage}[t]{0.9\linewidth}
\raggedright 1 – 100 μm
\end{minipage}
\end{enumerate}
\vspace{-0.3\baselineskip}
\begin{small}
\textit{คำตอบ: ก}
\end{small}

\item วัตถุ A มีอุณหภูมิ 3000 K และวัตถุ B มีอุณหภูมิ 6000 K อัตราส่วน λ\_max(A) / λ\_max(B) มีค่าเท่าใด?
\begin{enumerate}[label=\alph*)., itemsep=0.3\baselineskip, topsep=0.3\baselineskip]
\item ก\ 
\begin{minipage}[t]{0.9\linewidth}
\raggedright 1/2
\end{minipage}
\item ข\ 
\begin{minipage}[t]{0.9\linewidth}
\raggedright 1
\end{minipage}
\item ค\ 
\begin{minipage}[t]{0.9\linewidth}
\raggedright 2
\end{minipage}
\item ง\ 
\begin{minipage}[t]{0.9\linewidth}
\raggedright 4
\end{minipage}
\end{enumerate}
\vspace{-0.3\baselineskip}
\begin{small}
\textit{คำตอบ: ค}
\end{small}

\item คนตาบอดสีชนิดหนึ่งไม่สามารถแยกแยะสีแดงออกจากสีเขียวได้ เกิดจากความผิดปกติของเซลล์ประสาทชนิดใด?
\begin{enumerate}[label=\alph*)., itemsep=0.3\baselineskip, topsep=0.3\baselineskip]
\item ก\ 
\begin{minipage}[t]{0.9\linewidth}
\raggedright เซลล์แท่งทั้งหมด
\end{minipage}
\item ข\ 
\begin{minipage}[t]{0.9\linewidth}
\raggedright เซลล์กรวยรับแสงแดงและเขียว
\end{minipage}
\item ค\ 
\begin{minipage}[t]{0.9\linewidth}
\raggedright เซลล์กรวยรับแสงน้ำเงิน
\end{minipage}
\item ง\ 
\begin{minipage}[t]{0.9\linewidth}
\raggedright เส้นประสาทตา
\end{minipage}
\end{enumerate}
\vspace{-0.3\baselineskip}
\begin{small}
\textit{คำตอบ: ข}
\end{small}

\item จอมอนิเตอร์คอมพิวเตอร์แสดงสีได้โดยการผสมแสง และใช้ระบบสีแบบ RGB ถ้าต้องการแสดงสีฟ้า (Cyan) หน้าจอต้องเปิดพิกเซลสีใด?
\begin{enumerate}[label=\alph*)., itemsep=0.3\baselineskip, topsep=0.3\baselineskip]
\item ก\ 
\begin{minipage}[t]{0.9\linewidth}
\raggedright แดงเท่านั้น
\end{minipage}
\item ข\ 
\begin{minipage}[t]{0.9\linewidth}
\raggedright เขียวเท่านั้น
\end{minipage}
\item ค\ 
\begin{minipage}[t]{0.9\linewidth}
\raggedright น้ำเงินเท่านั้น
\end{minipage}
\item ง\ 
\begin{minipage}[t]{0.9\linewidth}
\raggedright เขียว + น้ำเงิน
\end{minipage}
\end{enumerate}
\vspace{-0.3\baselineskip}
\begin{small}
\textit{คำตอบ: ง}
\end{small}

\item ถ้าวัตถุร้อนจัดมี λ\_max = 500 nm เมื่อวัตถุเย็นลงจนมี λ\_max = 1000 nm อุณหภูมิจะลดลงกี่เท่า?
\begin{enumerate}[label=\alph*)., itemsep=0.3\baselineskip, topsep=0.3\baselineskip]
\item ก\ 
\begin{minipage}[t]{0.9\linewidth}
\raggedright 2 เท่า
\end{minipage}
\item ข\ 
\begin{minipage}[t]{0.9\linewidth}
\raggedright 4 เท่า
\end{minipage}
\item ค\ 
\begin{minipage}[t]{0.9\linewidth}
\raggedright 8 เท่า
\end{minipage}
\item ง\ 
\begin{minipage}[t]{0.9\linewidth}
\raggedright 16 เท่า
\end{minipage}
\end{enumerate}
\vspace{-0.3\baselineskip}
\begin{small}
\textit{คำตอบ: ก}
\end{small}

\item รังสีแกมมาที่ใช้ในทางการแพทย์มีความยาวคลื่นประมาณ 10⁻¹² m รังสีอินฟราเรดมีความยาวคลื่นประมาณ 10⁻⁵ m อัตราส่วน f\_gamma / f\_IR มีค่าเท่าใด?
\begin{enumerate}[label=\alph*)., itemsep=0.3\baselineskip, topsep=0.3\baselineskip]
\item ก\ 
\begin{minipage}[t]{0.9\linewidth}
\raggedright 10⁻³
\end{minipage}
\item ข\ 
\begin{minipage}[t]{0.9\linewidth}
\raggedright 10³
\end{minipage}
\item ค\ 
\begin{minipage}[t]{0.9\linewidth}
\raggedright 10⁷
\end{minipage}
\item ง\ 
\begin{minipage}[t]{0.9\linewidth}
\raggedright 10⁷
\end{minipage}
\end{enumerate}
\vspace{-0.3\baselineskip}
\begin{small}
\textit{คำตอบ: ง}
\end{small}

\item ช่วงความยาวคลื่นของแสงสีส้มอยู่ที่ประมาณเท่าใด?
\begin{enumerate}[label=\alph*)., itemsep=0.3\baselineskip, topsep=0.3\baselineskip]
\item ก\ 
\begin{minipage}[t]{0.9\linewidth}
\raggedright 400 – 500 nm
\end{minipage}
\item ข\ 
\begin{minipage}[t]{0.9\linewidth}
\raggedright 500 – 565 nm
\end{minipage}
\item ค\ 
\begin{minipage}[t]{0.9\linewidth}
\raggedright 565 – 590 nm
\end{minipage}
\item ง\ 
\begin{minipage}[t]{0.9\linewidth}
\raggedright 590 – 625 nm
\end{minipage}
\end{enumerate}
\vspace{-0.3\baselineskip}
\begin{small}
\textit{คำตอบ: ค}
\end{small}

\item ทำไมท้องฟ้าในเวลากลางวันจึงเห็นเป็นสีฟ้า?
\begin{enumerate}[label=\alph*)., itemsep=0.3\baselineskip, topsep=0.3\baselineskip]
\item ก\ 
\begin{minipage}[t]{0.9\linewidth}
\raggedright เพราะน้ำในบรรยากาศดูดกลืนแสงแดง
\end{minipage}
\item ข\ 
\begin{minipage}[t]{0.9\linewidth}
\raggedright เพราะอนุภาคในบรรยากาศกระเจิงแสงน้ำเงินมากกว่าแสงสีอื่น
\end{minipage}
\item ค\ 
\begin{minipage}[t]{0.9\linewidth}
\raggedright เพราะแสงแดงทะลุผ่านบรรยากาศได้ดีกว่า
\end{minipage}
\item ง\ 
\begin{minipage}[t]{0.9\linewidth}
\raggedright เพราะเมฆสะท้อนแสงน้ำเงิน
\end{minipage}
\end{enumerate}
\vspace{-0.3\baselineskip}
\begin{small}
\textit{คำตอบ: ข}
\end{small}

\item รังสีใดมีอำนาจทำลายเซลล์มากที่สุด?
\begin{enumerate}[label=\alph*)., itemsep=0.3\baselineskip, topsep=0.3\baselineskip]
\item ก\ 
\begin{minipage}[t]{0.9\linewidth}
\raggedright อินฟราเรด
\end{minipage}
\item ข\ 
\begin{minipage}[t]{0.9\linewidth}
\raggedright ไมโครเวฟ
\end{minipage}
\item ค\ 
\begin{minipage}[t]{0.9\linewidth}
\raggedright อัลตราไวโอเลต
\end{minipage}
\item ง\ 
\begin{minipage}[t]{0.9\linewidth}
\raggedright วิทยุ
\end{minipage}
\end{enumerate}
\vspace{-0.3\baselineskip}
\begin{small}
\textit{คำตอบ: ค}
\end{small}

\item หลอดไฟ LED สีม่วงเมื่อส่องลงบนกระดาษสีเหลืองจะเห็นกระดาษเป็นสีใด?
\begin{enumerate}[label=\alph*)., itemsep=0.3\baselineskip, topsep=0.3\baselineskip]
\item ก\ 
\begin{minipage}[t]{0.9\linewidth}
\raggedright แดง
\end{minipage}
\item ข\ 
\begin{minipage}[t]{0.9\linewidth}
\raggedright เขียว
\end{minipage}
\item ค\ 
\begin{minipage}[t]{0.9\linewidth}
\raggedright น้ำเงิน
\end{minipage}
\item ง\ 
\begin{minipage}[t]{0.9\linewidth}
\raggedright ดำหรือมืด
\end{minipage}
\end{enumerate}
\vspace{-0.3\baselineskip}
\begin{small}
\textit{คำตอบ: ง}
\end{small}

\item ดาวดวงหนึ่งมีอุณหภูมิ 12000 K ความยาวคลื่นที่ความเข้มสูงสุดจะอยู่ในช่วงสเปกตรัมใด? (กำหนด b = 2.898 × 10⁻³ m·K)
\begin{enumerate}[label=\alph*)., itemsep=0.3\baselineskip, topsep=0.3\baselineskip]
\item ก\ 
\begin{minipage}[t]{0.9\linewidth}
\raggedright อินฟราเรด
\end{minipage}
\item ข\ 
\begin{minipage}[t]{0.9\linewidth}
\raggedright แสงที่ตามองเห็น (ช่วงม่วง-น้ำเงิน)
\end{minipage}
\item ค\ 
\begin{minipage}[t]{0.9\linewidth}
\raggedright อัลตราไวโอเลต
\end{minipage}
\item ง\ 
\begin{minipage}[t]{0.9\linewidth}
\raggedright รังสีเอกซ์
\end{minipage}
\end{enumerate}
\vspace{-0.3\baselineskip}
\begin{small}
\textit{คำตอบ: ข}
\end{small}

\end{enumerate}
\end{document}